% ============================================================================
% preamble.tex — 包、字体、排版、代码环境、自定义彩框
% 《Airbnb 面试通关手册》
% ============================================================================

% --- 中文支持 (XeLaTeX) ---
\usepackage{ctex}
\usepackage{fontspec}

% 等宽字体：DejaVu Sans Mono（渲染代码 / 框图制表符；需安装到系统字体）
\setmonofont{DejaVu Sans Mono}[Scale=0.92]

% --- 页面布局 ---
\usepackage[
  a4paper,
  top=2.4cm, bottom=2.4cm,
  left=2.3cm, right=2.3cm
]{geometry}

% --- 数学 ---
\usepackage{amsmath, amssymb, amsthm}
\usepackage{mathtools}

% --- 图表 ---
\usepackage{graphicx}
\usepackage{float}
\usepackage{booktabs}
\usepackage{longtable}
\usepackage{tabularx}
\usepackage{array}
\usepackage{multirow}

% --- 绘图 ---
\usepackage{tikz}
\usetikzlibrary{positioning, calc, arrows.meta, fit, backgrounds, shapes.geometric}

% --- 颜色 ---
\usepackage{xcolor}
\definecolor{abred}{HTML}{FF5A5F}     % Airbnb 品牌红
\definecolor{abdark}{HTML}{484848}
\definecolor{abgray}{HTML}{767676}
\definecolor{codebg}{HTML}{F7F7F9}
\definecolor{codekw}{HTML}{0033B3}
\definecolor{codestr}{HTML}{067D17}
\definecolor{codecmt}{HTML}{8C8C8C}
\definecolor{codetype}{HTML}{8000A0}

% --- 超链接与交叉引用 ---
\usepackage{hyperref}
\hypersetup{
  colorlinks=true,
  linkcolor=abred!70!black,
  citecolor=green!40!black,
  urlcolor=cyan!55!black,
  bookmarksnumbered=true,
}
\usepackage[nameinlink]{cleveref}

% --- 代码高亮 (listings, 纯 LaTeX, 无需 python) ---
\usepackage{xstring}
\usepackage{listings}
\lstset{
  basicstyle=\ttfamily\small,
  backgroundcolor=\color{codebg},
  keywordstyle=\color{codekw}\bfseries,
  stringstyle=\color{codestr},
  commentstyle=\color{codecmt}\itshape,
  numberstyle=\tiny\color{abgray},
  numbers=left, numbersep=8pt,
  frame=single, framesep=5pt, rulecolor=\color{gray!30},
  breaklines=true, breakatwhitespace=false,
  showstringspaces=false,
  tabsize=2,
  columns=fullflexible,
  keepspaces=true,
  xleftmargin=14pt, xrightmargin=4pt,
  aboveskip=8pt, belowskip=8pt,
}
% C++ 别名（含常见 STL 类型着色）
\lstdefinestyle{cpp}{
  language=C++,
  morekeywords={auto, vector, string, unordered_map, unordered_set, map, set,
    priority_queue, queue, deque, pair, size_t, int64_t, nullptr, override,
    constexpr, string_view},
  emph={int, long, bool, char, double, void, struct, class},
  emphstyle=\color{codetype},
}
\lstnewenvironment{cpp}{\lstset{style=cpp}}{}

% --- 彩色定理框 ---
\usepackage[most]{tcolorbox}
\tcbuselibrary{breakable, skins}

% ============================================================================
% 自定义结构化彩框：本书面经的核心“设计思想”载体
% ============================================================================

% 设计思想 / 核心心法（红框）
\newtcolorbox{idea}[1][核心思想]{
  enhanced, breakable, arc=3pt,
  before skip=10pt, after skip=10pt,
  colback=abred!5, colframe=abred!80!black,
  boxrule=0.8pt, left=8pt, right=8pt, top=6pt, bottom=6pt,
  fonttitle=\bfseries, title={【设计思想】#1},
}

% 陷阱 / 易错点（橙框）
\newtcolorbox{pitfall}[1][常见陷阱]{
  enhanced, breakable, arc=3pt,
  before skip=10pt, after skip=10pt,
  colback=orange!6, colframe=orange!70!black,
  boxrule=0.7pt, left=8pt, right=8pt, top=6pt, bottom=6pt,
  fonttitle=\bfseries, title={【陷阱】#1},
}
% 追问 / Follow-up（蓝框）
\newtcolorbox{followup}[1][面试官追问]{
  enhanced, breakable, arc=3pt,
  before skip=10pt, after skip=10pt,
  colback=blue!4, colframe=blue!55!black,
  boxrule=0.7pt, left=8pt, right=8pt, top=6pt, bottom=6pt,
  fonttitle=\bfseries, title={【追问】#1},
}
% 复杂度 / 要点速记（绿框）
\newtcolorbox{keypoint}[1][要点]{
  enhanced, breakable, arc=3pt,
  before skip=10pt, after skip=10pt,
  colback=green!4, colframe=green!45!black,
  boxrule=0.7pt, left=8pt, right=8pt, top=6pt, bottom=6pt,
  fonttitle=\bfseries, title={【要点】#1},
}

% ============================================================================
% 题目标题命令：\problem{编号}{标题}{难度}
% ============================================================================
\newcommand{\difficulty}[1]{%
  \IfStrEqCase{#1}{%
    {Easy}{\textcolor{green!50!black}{\textbf{Easy}}}%
    {Medium}{\textcolor{orange!80!black}{\textbf{Medium}}}%
    {Hard}{\textcolor{abred!80!black}{\textbf{Hard}}}%
  }[\textbf{#1}]%
}
\newcommand{\problem}[3]{%
  \section{#2}%
  \noindent{\small\textcolor{abgray}{LeetCode #1 \quad|\quad 难度：}\difficulty{#3}}\par
  \vspace{4pt}\hrule height 0.4pt \vspace{8pt}%
}

% 常用记号
\newcommand{\bigO}[1]{\ensuremath{O(#1)}}
\newcommand{\code}[1]{\texttt{#1}}
